\begin{afterword}
\summaryhead

A short story. Jeffreyssai, a master of the beisutsukai, receives a guest from another Conspiracy, an old acquaintance whose counsel has always been good. She is a domain expert, she says, and certainty is not needed for her advice to count. She jokes that she has statistics, then admits that only he can decide, but says he should hold his ritual for seriously considering the abandonment of a long-held premise of his life. He concedes. She leaves with ``Farewell, sensei'' and ``It won't be the same.''

He cancels his classes and spends the day idle, letting his thoughts drift while blocking every thought he has had before. The next morning he goes down to a bare white room. On the wall are plaques of maxims and a tally: in seven earlier sessions he changed his mind twice. He sits facing the blank wall, rehearses mnemonics for his principles, techniques and past oversights, and asks himself the question. The story ends there.

\responsehead

The story is fiction, and it is judged here for what it shows, not for its lack of evidence. What it shows is the procedure that ``Crisis of Faith'' prescribed the day before: rest the day before, set aside quiet time, clear the mind of every thought already had, then ask. It shows that procedure faithfully. Two details fit the idea well: a guest who is not a rationalist moves a master by her record and her expertise, and the room's tally counts five kept beliefs as ordinary results.

That tally matters for the post before it. ``Crisis of Faith'' asked for a crisis that could ``go either way,'' and then counted Aumann's kept faith as a failure, because it took the faith to be false. The story makes plain what that example left unsaid. Most sessions end with the belief kept, and that is no failure; the sign of a missing ability would be never changing at all.

What the story leaves out is the part that would test it. The premise under review is not named, the reasoning is not shown, and the result is not given. The principles and techniques are counted (seven and sixty-two) but none is stated. So the story shows the preparation, and the reader cannot see whether the preparation leads anywhere. The preceding post described no real crisis either. This story is the nearest thing to one, and it stops at the first question.

Its place in the book adds to this. It closes the sequence ``Letting Go'' and the book \textsc{How to Actually Change Your Mind}. The book's last scene is a man about to consider changing his mind. It is also the reader's first meeting with Jeffreyssai and the beisutsukai, whose stories the author had begun earlier in 2008 but which the book places later, so the names arrive unexplained.

In short: a faithful dramatization of the Crisis of Faith procedure, clearer than that post that a kept belief is no failure, but ending before the one step that would show whether the procedure works.
\end{afterword}
